{\scriptsize
        \begin{tabularx}{\textwidth}{|p{0.25\columnwidth}|X|}
            \hline 
            \textbf{Notions et contenus} & \textbf{Capacités exigibles} \\
            \hline 
            \multicolumn{2}{|l|}{\textbf{\text { Ondes électromagnétiques dans les milieux matériels }}} \\
            \hline
            Propagation d'une onde électromagnétique plane harmonique unidirectionnelle dans un conducteur ohmique de conductivité réelle. Effet de peau dans un conducteur ohmique.  & 
            Identifier une analogie avec un phénomène de diffusion.
Établir la relation de dispersion des ondes électromagnétiques dans un conducteur ohmique à basses fréquences.
Associer l'atténuation de l'onde dans le milieu conducteur à une dissipation d'énergie. Estimer l'ordre de grandeur de l'épaisseur de peau du cuivre à différentes fréquences. \\
            \hline 
            Propagation d'une onde électromagnétique plane harmonique transverse et unidirectionnelle dans un plasma dilué. Conductivité électrique complexe.   & 
            Justifier la neutralité électrique locale du plasma en présence d'une onde transverse. Établir l'expression de la conductivité électrique complexe du plasma. Interpréter énergétiquement le caractère imaginaire pur de la conductivité électrique complexe du plasma. \\
        \hline 
        Relation de dispersion. Pulsation plasma. Domaine de transparence. Domaine réactif, onde évanescente. & 
        Établir la relation de dispersion des ondes planes progressives harmoniques transverses.
        Exprimer la vitesse de phase et la vitesse de groupe d'un paquet d'ondes dans le domaine de transparence du plasma. Interpréter la pulsation plasma comme une pulsation de coupure.
        Citer les caractéristiques d'une onde stationnaire évanescente. Justifier que, dans le domaine réactif, une onde électromagnétique harmonique ne transporte aucune puissance en moyenne. \\ 
        \hline
        \end{tabularx}
\vspace{0.5cm}
}
